\documentclass[10pt,a4paper]{amsart}
\usepackage[margin=19mm]{geometry}
\usepackage{amsmath,amssymb,mathtools}
\usepackage[scr=rsfs,cal=euler]{mathalfa}
\usepackage{xfrac}
\usepackage[unicode,hidelinks,bookmarks=false]{hyperref}
\newcommand{\ie}{i.e.\ }
\newcommand{\cf}{cf.\ }
\newcommand{\resp}{respectively\ }
\newcommand{\Subsection}{\subsection}
\providecommand{\nobreakdash}{\nobreak\hbox{-}}
\setlength{\emergencystretch}{2em}
\multlinegap0pt
\usepackage{bm}
\usepackage{bbm}
\usepackage{dsfont}
\usepackage{xspace}
\usepackage{cancel}
\usepackage{mathtools}
\usepackage{amsthm}
\makeatletter
\@ifundefined{lemma}{\newtheorem{lemma}{Lemma}}{}
\makeatother
\title[Kac's Lemma and countable generators for actions of countabl]{Kac's Lemma and countable generators for actions of countable groups}
\author{Tom Meyerovitch, Benjamin Weiss}
\date{Source: arXiv:2410.18488v2 (CC BY 4.0)}
\begin{document}
\maketitle

\noindent\emph{Source and licence.} Kac's Lemma and countable generators for actions of countable groups. Version arXiv:2410.18488v2 (CC BY 4.0), distributed under the Creative Commons Attribution 4.0 International licence, as recorded in the arXiv licence metadata for this exact version and shown on the abstract page https://arxiv.org/abs/2410.18488v2. Full rights notice: licenses/arxiv-2410.18488-CC-BY-4.0.txt.\par\smallskip

\noindent\textit{Editorial scope.} Counted unit: the Lemma whose source label is \texttt{lem:sweep\_out\_partition} in the article (source0.tex lines 592--596, source label \texttt{lem:sweep\_out\_partition}), delivered together with the article's own complete original proof of it (source0.tex lines 597--640, one \texttt{proof} environment of 44 lines). No mathematics is written, repaired or re-derived: statement and proof are the article's own text line for line. Required context retained uncounted: none. Every definition, display and label of those lines is retained unchanged. Omitted from this package and not redistributed: 1--591, 641--833. Nothing inside a retained excerpt is omitted or altered: every retained body is delivered exactly as the article's lines are written, so no omission receipt of any kind accompanies this package. Citations: the retained excerpts carry no citation command at all, so no bibliography entry of the article is transcribed and no reference is redistributed. Reference adaptation: none was needed and none was made; every reference inside the delivered unit resolves inside it, so no unresolved reference or citation is produced. Printed slips kept literal: no printed word, symbol, hypothesis or number of the counted statement or of its proof was altered. Layout and class: the article's own class is \texttt{amsart} with packages including amsmath, amsthm, amsopn, amssymb, microtype, mathrsfs, mathscinet, bbm, enumerate, tikz, etoolbox, intcalc, geometry, caption; the wrapper is the lane's pinned \texttt{amsart} editorial wrapper at 10pt on A4, which restates the theorem-like environments the retained text needs and adds the article's own macro definitions that the retained text needs (none), with extra wrapper packages (bm, bbm, dsfont, xspace, cancel, mathtools). The delivered numbering is the unit's own. Assets: no retained span contains a figure environment, a graphics-inclusion command, a PDF-page inclusion, a table or a TikZ picture; no image file is copied into this package, the package declares no asset dependency and no image file is redistributed with it. Licence: version arXiv:2410.18488v2 of this article is distributed under the Creative Commons Attribution 4.0 International licence, as recorded in the arXiv licence metadata for this exact version and shown on the abstract page https://arxiv.org/abs/2410.18488v2. A full rights notice is delivered at licenses/arxiv-2410.18488-CC-BY-4.0.txt. Characters: every retained excerpt is pure ASCII, so the delivered engine, XeLaTeX with its default Latin Modern text font, reports no missing character.
\begin{lemma}\label{lem:sweep_out_partition}
    Let $(X,\mathcal{B},\mu,T)$ be a probability preserving action of the countable group $\Gamma$ on a standard  probability space $(X,\mathcal{B},\mu)$. Assume that almost every point in $X$ has an infinite orbit. 
    Then for any $\epsilon >0$ there exists a countable collection of pairwise disjoint sweep-out subsets of  $X$, each of which has measure at most $\epsilon$.
    %= \biguplus_{n=1}A_n$ so that  each $A_n$ is a sweep-out set with $\mu(A_n) \le \epsilon$.
\end{lemma}

\begin{proof}
    Since $(X,\mathcal{B},\mu)$ is a standard probability space, we can assume by  Carath{\'e}odory's theorem that $X=[0,1]$,  that $T$ is a $\Gamma$-action on $[0,1]$ by Borel automorphisms that $\mu$ is a  $\Gamma$-invariant Borel-measure on $[0,1]$ and that $\mathcal{B}$ is the $\mu$-completion of the Borel $\sigma$-algebra on $[0,1]$.
    Let $\mathcal{I} \subseteq \mathcal{B}$ denote the $\sigma$-algebra of $T$-invariant sets. Given $A \in \mathcal{B}$, the conditional probability $\mu(A \mid \mathcal{I})(x)$ is defined $\mu$-almost everywhere, and $x \mapsto \mu(A \mid \mathcal{I})(x)$ is a $\mu$-measurable  almost everywhere defined function.

Fix $0 < \epsilon < 1$.

Let 
\[
f_1(x) = \inf \{ t >0~:~ \mu( [0,t] \mid \mathcal{I})(x) \ge \epsilon/2\}.
\]
The almost-everywhere defined function $f_1:X \to [0,1]$ is $\mathcal{I}$-measurable, thus $\mathcal{B}$-measurable.

Since almost every point in $X$ has an infinite orbit,  we conclude that almost surely the conditional probability $\mu(\cdot \mid \mathcal{I})(x)$ has no atoms, and so  for almost every $x \in X$ the function $t \mapsto  \mu([0,t] \mid \mathcal{I})(x)$ is continuous. Thus $\mu([0,f_1(x)] \mid \mathcal{I})(x) = \epsilon/2$ for $\mu$-almost every $x \in X$. In particular, since $\epsilon /2 < 1$, it follows that $f_1(x) < 1$ almost everywhere.
Similarly, by induction  we can define for every $n \ge 2$ an almost everywhere defined measurable function  $f_{n}:X \to [0,1]$ by
\[
f_n(x) = \inf \{ t >f_{n-1}(x) ~:~ \mu( [f_{n-1}(x),t) \mid \mathcal{I})(x) \ge \epsilon/2^n\}.
\]
Similarly to the case $n=1$, using that $\mu(
\cdot 
\mid \mathcal{I})$ is nonatomic almost surely, we have that  $\mu$-almost everywhere
\[ 0 < f_1(x) < f_2(x) < \ldots < f_n(x) < \ldots < 1,\]
and also for  $\mu$-almost every $x \in X$ we have
\[ \mu\left( (f_{n-1}(x),f_n(x)] \mid \mathcal{I}\right)(x) = \frac{\epsilon}{2^n}. \]
Let
\[
A_1 = \{ x \in [0,1]~:~ x \le f_1(x)\}
\]
and for every $n \ge 2$
\[
A_n = \{ x \in [0,1]~:~ f_{n-1}(x) < x \le f_n(x)\}.
\]
Then $A_1,\ldots,A_n,\ldots$ are pairwise disjoint $\mathcal{B}$-measurable sets with $\mu(A_n) = \frac{\epsilon}{2^n}$.
It remains to check that each $A_n$ is a sweep-out set.

Indeed, since $\bigcup_{\gamma \in \Gamma}T_\gamma(A_n) \in \mathcal{I}$ it follows that $\mu$-almost everywhere we have

$\mu( \bigcup_{\gamma \in \Gamma}T_\gamma(A_n) \mid \mathcal{I}) (x)\in \{0,1\}$.

Since $\mu( \bigcup_{\gamma \in \Gamma}T_\gamma(A_n) \mid \mathcal{I}) (x) \ge \epsilon/2^n$ almost-everywhere, we conclude that indeed
 each $A_n$ is a sweep-out set.



\end{proof}

\end{document}
